\section{Discussion}
\label{sec:discussion}

\paragraph{What a sensing target determines.}
The method is useful when expected sensing usage is the intended design requirement. A target in observation count or cost gives an engineer an interpretable input, and scale equivariance makes the corresponding transformation of $w$ predictable when reward units change. The target does not determine which uncertainty should be reduced or maximize reward under a cap. Ordinary information gain may reward irrelevant sensing. In the distractor study, scoring information on reward-equivalence classes removes that spending at every swept weight, while the budget controls how much sensing is bought overall (Section~\ref{sec:distractor}).

Even within the chosen family, calibration is not a general policy optimizer. A plateau can contain many weights with the same usage, and the selected endpoint mixture need not maximize reward among all mixtures meeting the target. The held-out budget study quantifies both distinctions, against a sampled constrained reference and against alternative mixtures over the calibration grid. Its targets come from a predeclared rule applied to that grid, not from an application-specified inspection requirement. Deployment would additionally require checking the supplied model and the stability of calibration under the intended operating conditions.

\paragraph{How to interpret the information-unit default.}
The recursive EFE identity fixes $w=1$ once reward and information units have been declared. It does not eliminate that convention. The weight belongs to the level set of its implicit budget $U(1)$, and inversion need not return the same weight. For the specified recursion, the formal equivalence is exact under its assumptions. The claim that $w=1$ is a useful default is empirical: it lies in the reward-maximizing tied bracket on three of five swept environments and is beaten on the other two. Its near-optimality in a random two-state family also depends on horizon and the chosen reward-gap criterion. Even at $H=3$, approximately $21\%$ of that family fails the study's criterion (Appendix~\ref{app:nearopt_horizon}).

The comparison between reward-tuned and success-tuned weights reflects different objectives. A larger weight can improve success while spending more on tests and reducing reward. Moderate transferred weights perform well on parts of the suite under shared reward units, but transfer at the top of the grid can be costly (Table~\ref{tab:transfer}). These observations support checking reward, success, and usage separately. They do not establish a universal operating point.

\paragraph{Which problem structures benefit.}
In the tested regimes, EFE helps when the available search can exploit informative observations that a reward-only planner undervalues. Multiple observation actions create opportunities to select useful tests, but neither their number nor state-space size alone predicts an advantage. On the largest Tileworld grid, the fixed-horizon reward-only planner takes no scans, making the separation a joint horizon-and-scale effect. On RockSample[11,11], the advantage disappears, and on the two larger RockSample instances an unbounded-horizon heuristic beats the tree-search agents. Navigation supplies a further negative case: no epistemic agent significantly beats the reward-only greedy mover at any tested grid size (Appendix~\ref{app:navigation}).

Planning depth and discounting also change the value of information. At one-step depth the recursive objective reduces to myopic reward plus information gain at $w=1$. The observed discount threshold in the benchmark sweep is not a general regime boundary. Appendix~\ref{app:interpretation} retains the detailed horizon, transfer, and discount comparisons. Model accuracy matters independently. EFE and reward-only Planning respond identically to every tested Tiger sensor mismatch, so that study supplies no EFE-specific robustness benefit. Overestimating sensor accuracy can trigger premature commits (Appendix~\ref{app:misspec}).

\paragraph{What the tree-search controls establish.}
Exact belief-tree enumeration limits practical planning depth, motivating MCTS-EFE. The principal comparison with POMCP matches simulation counts, with different costs per simulation. Tuning POMCP's exploration constant recovers most of the initial Tiger gap. The narrower comparison after selecting both solvers' constants uses disjoint tuning seeds, but its evaluation runs had already been inspected, so the Tiger result is retrospective evidence. On Diagnosis and Tileworld, the gap survives the tested constants and information-gain rollouts. The ablation finds no detectable effect of either in-tree information gain or max-backup at five seeds after multiplicity correction.

The remaining components, EFE leaf evaluation, exact beliefs, and exact commit-child values, were held fixed together. The experiments therefore evaluate their joint configuration and do not isolate directed observation selection as the cause of the advantage. Nor do they establish superiority over POMDP solvers in general. Full solver configurations, timings, controls, and outcomes appear in Appendices~\ref{app:interpretation} and~\ref{app:pomcp}.

\paragraph{Structural and practical limits.}
The equivalence covers observation and navigation actions that preserve the hidden state. When sensing changes that state, the epistemic term must use transition-aware predictive and posterior distributions. The transition--observation coupling diagnostic $\Delta_T$ identifies a failure of the reduction's hypothesis, but has not been proved to be an additive correction or a decision-error bound. Small transition probabilities do not guarantee a small diagnostic: with perfect sensing, a posterior support change can make the divergence infinite. Example~\ref{ex:destructive} gives a concrete failure, and Appendix~\ref{app:theory_details} develops the full scope discussion.

All experiments assume a known generative model. The large Structural Inspection state space demonstrates feasibility on a synthetic benchmark with a hand-coded leaf rule, not on a deployed inspection system. Model learning, more efficient search with many sensing actions, transition-aware epistemic objectives, and validation against an actual operational sensing requirement remain open work. These limits distinguish the exact mathematical results from the empirical guidance drawn from this suite.

\section{Conclusion}

An expected sensing-usage requirement can calibrate the information-gain weight of a $\rho$-POMDP policy family. The operational shadow price is a crossing threshold, estimated by a grid bracket. Under stated conditions, mixing the bracket's endpoint policies meets a usage-gap target in expectation. The receding-horizon planner is scale equivariant, and a projected decaying-step controller has a stationary-convergence guarantee under explicit hypotheses. Re-adaptation after reward-scale shifts is an empirical result, including a useful reset heuristic that the theorem does not cover.

The held-out study shows accurate usage calibration but also its reward cost. Meeting an equality target can earn less than spending below it, and the selected mixture need not be optimal within the family or against the constrained reference. Calibration is appropriate when expected usage itself is the requirement.

The recursive EFE objective supplies an information-unit anchor at $w=1$, with exact reward--information identification under log scoring and an explicit convention for ordinary rewards. This weight is competitive in several benchmarks and suboptimal in others. Its practical value depends on the observation structure, planning depth, and model assumptions. The resulting design procedure is to state the intended usage, calibrate within a specified policy family, and assess usage, reward, and success together.
